\section{The Clebsch cubic bridge}
\label{sec:cubic}

Let \(G\) be the order-\(60\) stabilizer of the \(H_3\)-invariant matching
in \(\operatorname{PGL}_2(\F_{11})\), and let \(\Omega\) be its
\(22\)-point matching orbit.  For a reference matching \(M_0\), put
\[
 W=\left\langle q_M-q_{M_0}:M\in\Omega\right\rangle .
\]
The two \(\PSL_2(\F_{11})\)-sheets define signs
\(\epsilon(M)\in\{\pm1\}\) and the symmetric tensor
\[
 T=\sum_{M\in\Omega}\epsilon(M)[q_M-q_{M_0}]^{\otimes3}
 \in\operatorname{Sym}^3(W).
\]

\begin{theorem}[Finite tensor bridge]
\label{thm:finite-tensor-bridge}
\claimid{HC-3}.
There is a \(G\)-equivariant isomorphism
\[
 S:\F_{11}^{\binom{[5]}2}\longrightarrow W
\]
such that, after using the standard pairing to identify the permutation
module with its dual and restricting along
\[
 E(y)_{\{i,j\}}=y_i+y_j,\qquad \sum_i y_i=0,
\]
one has
\[
 E^*(S^{-1})^{\otimes3}T=4\sigma_3.
\]
In particular, the signed matching cubic and the Clebsch cubic span the
same nonzero invariant line over \(\F_{11}\).
\end{theorem}

\begin{proof}
Conjugation on the normalizers of the five Klein four subgroups identifies
\(G\) with \(A_5\) in its degree-five action.  Hence it acts on the ten
two-subsets.  The matching construction gives a second ten-dimensional
representation on \(W\).  On the classes
\(1,2A,3A,5A,5B\), both characters are
\[
 (10,2,1,0,0),
\]
so both modules decompose as \(\mathbf1\oplus V_4\oplus V_5\).
This proves that an equivariant isomorphism exists.  To fix its three
irreducible scalings, we solve
\(\rho_W(g)S=S\rho_P(g)\).  The recorded invertible solution satisfies
this identity for all \(60\) elements.

The transported restriction of \(T\) is an \(A_5\)-invariant cubic on
\(V_4\).  Such cubics form one line: the degree-three symmetric orbit sums
on five variables have exponent types \(3\), \(2+1\), and \(1+1+1\), and
after imposing \(\sum y_i=0\) only the line generated by
\(\sigma_3\) remains.  It therefore suffices to determine one scalar.
One mixed polarized entry of the transported tensor is \(6\), while the
corresponding entry of \(\sigma_3\) is \(-1/3=7\).  Since
\(4\cdot7=6\) in \(\F_{11}\), the scalar is \(4\).

The displayed character argument and one-entry calculation are the proof
mechanism.  A primary exact calculation and an independent reconstruction
audit the chosen intertwiner on all \(60\) elements and all \(64\)
contracted entries.  Lean optionally checks the final literal tensor
equality; no formal result is used for the representation-theoretic
reduction.
\end{proof}

\begin{remark}[Normalization]
The characteristic-zero Gaunt scalar in
Theorem~\ref{thm:harmonic-clebsch} has denominator
divisible by \(11\).  The theorem therefore identifies a common integral
cubic line; it is not obtained by reducing that rational scalar modulo
\(11\).
\end{remark}
